% Section experiencia corta - DATA SCIENCE / ANÁLISIS DE DATOS
\sectionTitle{Experiencias Profesionales}{\faSuitcase}

\begin{experiences}

%\experience
%{Agosto 2025} {Data Analyst}{Dirección de Información Agraria y Estudios Económicos}{DRA-Ayacucho}
%{Abril 2025} {
%	Análisis de datos económicos y estadísticos del sector agrario regional para generación de información estratégica.
%	\begin{itemize}
%		\item Procesamiento, limpieza y validación de bases de datos estructuradas del sector agrario (30,000+ registros)
%		\item Desarrollo de scripts en R para automatización de procesos de ETL y generación de reportes
%		\item Elaboración de boletines estadísticos mensuales con análisis de coyuntura económica regional
%		\item Análisis cuantitativo de indicadores: producción agrícola, precios en chacra, empleo rural
%		\item Creación de visualizaciones con ggplot2 para comunicación de hallazgos a tomadores de decisión
%		\item Colaboración con equipos técnicos para consolidación de información económica intersectorial
%	\end{itemize}
%}
%{\footnotesize{\emph{Stack técnico:} R (dplyr, tidyr, ggplot2), Excel avanzado, Stata, SQL básico}}

\experience
{Ahora} {Analista de Datos Freelance}{Proyectos Personales}{Remoto}
{Enero 2024} {
	Desarrollo de proyectos de análisis de datos y visualización para aprendizaje continuo y portafolio técnico.
	\begin{itemize}
		\item Dashboard interactivo de coyuntura económica regional con datos del INEI/BCRP usando Shiny
		\item Web scraping de datos económicos públicos y automatización de actualización de datasets
		\item Análisis exploratorio de datos de encuestas nacionales (ENAHO, ENE) con R y Python
		\item Creación de visualizaciones interactivas y storytelling con datos económicos
		\item Documentación de proyectos en GitHub con README detallados y notebooks reproducibles
	\end{itemize}
}
{\footnotesize{\emph{Tecnologías:} R, Python (Pandas, Matplotlib, Seaborn), SQL, Git, Shiny, Streamlit}}
\emptySeparator

\experience
{Julio 2024} {Docente con Componente Analítico}{Universidad Nacional Autónoma de Huanta}{Huanta}
{Abril 2024} {
	Docencia aplicada con fuerte componente de análisis de datos y visualización de información cuantitativa.
	\begin{itemize}
		\item Enseñanza de análisis de datos para gestión de proyectos con Excel avanzado
		\item Diseño de ejercicios prácticos de análisis exploratorio y presentación de resultados
		\item Elaboración de dashboards y reportes visuales para seguimiento de indicadores de proyectos
		\item Capacitación en herramientas de organización y visualización de información estructurada
	\end{itemize}
}
{\footnotesize{\emph{Herramientas enseñadas:} Excel avanzado (tablas dinámicas, fórmulas complejas, gráficos), Google Sheets}}
\emptySeparator

%%%%%%%%%%%%%%%%%%%%%%%%%%%%%%%%%%%%%%%%
% EXPERIENCIAS ADICIONALES - Descomentar según relevancia
%%%%%%%%%%%%%%%%%%%%%%%%%%%%%%%%%%%%%%%%

%\experience
%{Marzo 2024} {Data Science Intern}{Startup Tecnológica}{Remoto}
%{Diciembre 2023} {
%	Prácticas en equipo de datos desarrollando modelos y análisis para producto digital.
%	\begin{itemize}
%		\item Desarrollo de modelos de clasificación con scikit-learn para segmentación de usuarios
%		\item Análisis de comportamiento de usuarios mediante análisis de cohortes y funnels de conversión
%		\item Creación de reportes automatizados con Python (reportlab) y visualizaciones con Plotly
%		\item Implementación de pruebas A/B y análisis estadístico de significancia
%		\item Contribución a documentación técnica y presentación de hallazgos en sprint reviews
%	\end{itemize}
%}
%{\footnotesize{\emph{Stack:} Python (Pandas, scikit-learn, Plotly), SQL (PostgreSQL), Git, Jupyter}}

%\experience
%{Noviembre 2023} {Asistente de Investigación Cuantitativa}{Centro de Investigación}{Ayacucho}
%{Mayo 2023} {
%	Apoyo en procesamiento y análisis de datos para proyectos de investigación económica aplicada.
%	\begin{itemize}
%		\item Limpieza y preparación de bases de datos de encuestas (ENAHO, ENAPRES) con Stata/R
%		\item Análisis estadístico descriptivo y elaboración de tablas para reportes de investigación
%		\item Programación de rutinas de análisis reproducibles con código documentado
%		\item Elaboración de gráficos publicables con ggplot2 siguiendo estándares académicos
%		\item Apoyo en documentación de metodologías y fuentes de datos
%	\end{itemize}
%}
%{\footnotesize{\emph{Software:} R (tidyverse, survey), Stata, Excel, LaTeX para documentación}}

%\experience
%{Agosto 2023} {Analista Junior de Business Intelligence}{Empresa Comercial}{Ayacucho}
%{Febrero 2023} {
%	Análisis de datos comerciales y desarrollo de reportes para seguimiento de ventas y operaciones.
%	\begin{itemize}
%		\item Extracción de datos de sistema ERP mediante consultas SQL para análisis comercial
%		\item Desarrollo de dashboards en Power BI para seguimiento de KPIs de ventas y rotación
%		\item Análisis de tendencias de ventas, estacionalidad y comportamiento de clientes
%		\item Elaboración de reportes automatizados para gerencia y áreas comerciales
%		\item Identificación de oportunidades de mejora en gestión de inventarios mediante análisis de datos
%	\end{itemize}
%}
%{\footnotesize{\emph{Herramientas:} Power BI, SQL Server, Excel avanzado, Python (básico para ETL)}}

%\experience
%{Enero 2023} {Practicante de Análisis de Datos}{Organismo Público}{Ayacucho}
%{Julio 2022} {
%	Prácticas preprofesionales en área de estadística e informática de entidad pública.
%	\begin{itemize}
%		\item Procesamiento de bases de datos administrativas para elaboración de anuarios estadísticos
%		\item Validación y consistencia de información entre diferentes sistemas institucionales
%		\item Elaboración de cuadros estadísticos y gráficos para publicaciones institucionales
%		\item Apoyo en desarrollo de indicadores de gestión y tableros de control
%		\item Documentación de procesos de manejo de información y diccionarios de datos
%	\end{itemize}
%}
%{\footnotesize{\emph{Software utilizado:} Excel, SPSS, Access, herramientas de visualización básicas}}

%\experience
%{Junio 2022} {Voluntario en Análisis de Datos}{ONG de Desarrollo}{Ayacucho}
%{Marzo 2022} {
%	Voluntariado técnico en análisis de datos de programas sociales para evaluación de impacto.
%	\begin{itemize}
%		\item Análisis de línea de base de programa de desarrollo productivo rural
%		\item Procesamiento de encuestas socioeconómicas y construcción de indicadores
%		\item Elaboración de perfiles estadísticos de beneficiarios
%		\item Apoyo en diseño de sistema de monitoreo con indicadores cuantitativos
%		\item Presentación de hallazgos a equipo técnico y stakeholders del proyecto
%	\end{itemize}
%}
%{\footnotesize{\emph{Herramientas:} R, Excel, visualizaciones para audiencias no técnicas}}

%\experience
%{Diciembre 2021} {Asistente de Laboratorio de Econometría}{Universidad}{Ayacucho}
%{Agosto 2021} {
%	Apoyo docente en curso de Econometría con énfasis en práctica con software estadístico.
%	\begin{itemize}
%		\item Asistencia a estudiantes en uso de Stata para estimación de modelos econométricos
%		\item Desarrollo de ejercicios prácticos y tutoriales de análisis de regresión
%		\item Revisión de trabajos de estudiantes: código, interpretación de resultados
%		\item Elaboración de material didáctico: do-files comentados, bases de datos de práctica
%	\end{itemize}
%}
%{\footnotesize{\emph{Temas cubiertos:} MCO, heterocedasticidad, series de tiempo, datos panel}}

\end{experiences}
